% ============================================================================
% Chapter 39: Online Safety and Digital Citizenship
% ============================================================================
\chapter{Online Safety and Digital Citizenship}
\label{ch:safety}

\begin{objectives}
  \item Understand the concept of digital citizenship
  \item Identify online threats: cyberbullying, scams, inappropriate content
  \item Protect personal information online
  \item Follow rules for responsible internet use
\end{objectives}

\bytetip[warn]{The internet is amazing, but it can also be dangerous if you're not careful. Being a good \textbf{digital citizen} means being safe, kind, and smart online.}

\section{What Is Digital Citizenship?}
\label{sec:digital-citizen}
\index{digital citizenship}

\begin{theorybox}[Digital Citizenship]
\textbf{Digital Citizenship}\index{digital citizenship} means using technology responsibly, ethically, and safely. A good digital citizen:
\begin{itemize}[leftmargin=*, label={\textcolor{LabGreen}{\ding{51}}}, itemsep=2pt]
  \item Treats others with \textbf{respect} online
  \item Protects their \textbf{personal information}
  \item Thinks before they \textbf{post or share}
  \item Reports \textbf{bullying or harmful content}
  \item Uses technology to \textbf{learn and create}, not harm
\end{itemize}
\end{theorybox}

\section{Staying Safe Online}
\label{sec:stay-safe}
\index{online safety}

% --- Figure 39.1: Online Safety Rules Cards ---
\begin{figure}[H]
\centering
\begin{tabularx}{\textwidth}{@{}*{4}{>{\centering\arraybackslash}X}@{}}
  \begin{tcolorbox}[colback=FunCoral!10, colframe=FunCoral, width=3.2cm, height=2.8cm, valign=center, halign=center, arc=8pt]
    {\Large\textcolor{FunCoral}{\faUserSecret}}\vspace{4pt}\\
    \textbf{Never share}\\personal info
  \end{tcolorbox}
  &
  \begin{tcolorbox}[colback=FunSky!10, colframe=FunSky, width=3.2cm, height=2.8cm, valign=center, halign=center, arc=8pt]
    {\Large\textcolor{FunSky}{\faLock}}\vspace{4pt}\\
    \textbf{Use strong}\\passwords
  \end{tcolorbox}
  &
  \begin{tcolorbox}[colback=LabGreen!10, colframe=LabGreen, width=3.2cm, height=2.8cm, valign=center, halign=center, arc=8pt]
    {\Large\textcolor{LabGreen}{\faUserFriends}}\vspace{4pt}\\
    \textbf{Only accept}\\real friends
  \end{tcolorbox}
  &
  \begin{tcolorbox}[colback=FunPurple!10, colframe=FunPurple, width=3.2cm, height=2.8cm, valign=center, halign=center, arc=8pt]
    {\Large\textcolor{FunPurple}{\faFlag}}\vspace{4pt}\\
    \textbf{Report}\\cyberbullying
  \end{tcolorbox}
\end{tabularx}
\caption{Four golden rules for staying safe online}
\label{fig:safety-rules}
\end{figure}

\section{Cyberbullying}
\label{sec:cyberbullying}
\index{cyberbullying}

\begin{warnbox}
\textbf{Cyberbullying}\index{cyberbullying} is using technology to harass, threaten, or humiliate someone. This includes:
\begin{itemize}[leftmargin=*, itemsep=2pt]
  \item Sending mean messages or comments
  \item Sharing someone's photos without permission
  \item Excluding someone from online groups to hurt them
  \item Creating fake profiles to impersonate someone
\end{itemize}
If you see cyberbullying: \textbf{don't respond}, \textbf{save evidence}, and \textbf{tell a trusted adult immediately}.
\end{warnbox}

\bytetip[tip]{The \textbf{Golden Rule} of the internet: Don't say or do anything online that you wouldn't say or do face-to-face. If you wouldn't show it to your grandmother, don't post it!}

\begin{funfact}
In many countries, cyberbullying is a \textbf{criminal offence}. In India, Section 67 of the IT Act can lead to up to \textbf{3 years in jail} for sharing offensive or harmful content online. The internet is NOT a lawless place!
\end{funfact}

\begin{reviewqs}
  \item What is digital citizenship? List three qualities of a good digital citizen.
  \item Name four types of personal information you should never share online.
  \item What is cyberbullying, and how should you respond if you see it?
  \item Why should you think twice before posting something online?
  \item What are three rules for staying safe on social media?
\end{reviewqs}

\begin{challenge}
Create a ``Digital Safety Pledge'' poster with at least 8 rules for safe internet use. Design it colourfully and hang it near your computer at home or school!
\end{challenge}
